\section{Exact transmission on the regular harmonic subspaces}
\label{sec:transmission}

The link, selection, and control results are now proved. Fix finite spectral
data as specified at the end of Section~\ref{sec:counting};
Lemma~\ref{lem:spectral-surplus} supplies such data. Thus the real spaces
\(V_l\), the finite range \(2\leq l\leq L\), and the selected
multiplicities \(M_l\) are fixed.
Keep the ordered dwell eigenbases and signed cycles \(\Pi_l\) chosen
in Section~\ref{sec:control}. Every matrix in this section acts in those
fixed round orthonormal coordinates.

Our task is to make the harmonic equation realize these cycles after
an entire period, including its long nonround dwell. First we choose the
pulse family using Corollary~\ref{prop:word}. Proposition~\ref{prop:metric}
then supplies the metric for every admitted sequence once the start is
sufficiently late. The remainder of this section proves uniform estimates
for those provisional metrics before choosing the actual sequence.

\subsection{The ideal pulse and its geometric realization}
We now choose the constants in the control proposition to agree with
the response of the radial equation.  Let \(p_\infty=m-1\), and
for \(2\leq l\leq L\) define
\begin{equation}
 \theta_l>0,\qquad
 \theta_l^2+p_\infty\theta_l=a_*^{-2}B_l,\qquad
 \kappa_l=\frac{a_*^{-2}}{p_\infty+2\theta_l}>0.
 \label{eq:round-root}
\end{equation}
Apply Corollary~\ref{prop:word} to the target \(\Pi\) with these \(\kappa_l\).
Restrict its parameter domain to a closed ball \(K\) centered at
zero and contained in \(U\).  Choose open intervals
\(I_1,\ldots,I_R\subset(0,1)\), in chronological order and with
disjoint closures, and smooth functions \(\beta_\nu\) compactly
supported in \(I_\nu\) with
\(\int_0^1\beta_\nu(\tau)\,d\tau=1\).  Put
\begin{align}
 z(\tau,h)&=\sum_{\nu=1}^{R}c_\nu(h)\beta_\nu(\tau),\nonumber\\
 R_l(\tau,h)&=\kappa_l\sum_{\nu=1}^{R}
   c_\nu(h)\beta_\nu(\tau)
   \left(\frac{B_l}{m}I+D_{J_\nu}^2|_{V_l}\right).
 \label{eq:pulse-profile}
\end{align}
These functions, with all derivatives that occur below, are bounded
on the fixed compact parameter set.  Let \(H_l\) solve
\begin{equation}
 \partial_\tau H_l=R_l(\tau,h)H_l,\qquad H_l(0,h)=I.
 \label{eq:ideal-pulse}
\end{equation}
For a matrix \(C\) on \(V_l\) with positive determinant, define
\[
 \mathcal N_l(C)=(\det C)^{-1/N_l}C.
\]
Each bump in \eqref{eq:pulse-profile} is a scalar function of time
times one constant matrix.  It therefore contributes
\(\exp\bigl(c_\nu(h)\kappa_l(B_lI/m+D_{J_\nu}^2)\bigr)\) to
\(H_l(1,h)\), with later bumps multiplying on the left.  For
positive-determinant matrices,
\(\mathcal N_l(C_2C_1)=\mathcal N_l(C_2)\mathcal N_l(C_1)\);
also \(\mathcal N_l(e^C)=e^{C_{\mathrm{tf}}}\) for every real
matrix \(C\).  Thus removing the determinant removes the scalar
part of each exponent.  Consequently
\begin{equation}
 \begin{gathered}
 \Phi(h):=\bigl(\mathcal N_l(H_l(1,h))\bigr)_{l=2}^{L}
          =\mathcal W(h),\\
 \Phi(0)=\Pi,\qquad D\Phi(0)\text{ is an isomorphism}.
 \end{gathered}
 \label{eq:ideal-word}
\end{equation}
The determinants of the \(H_l\) are positive, since they are
exponentials of integrals of real traces.


Apply Proposition~\ref{prop:metric} to this fixed pulse family. Its proof
supplies a threshold \(j_{\rm geom}\) such that every \(j_0\geq j_{\rm geom}\)
and every parameter sequence in \(K\) give a smooth complete metric with
nonnegative Ricci curvature. Use exactly the radial scale \(a\), cutoff,
and five-part schedule \eqref{eq:periods} fixed there. In particular,
\(d_g(0,x)=r(x)\), \(b=a'/a\geq-1/4\), and the angular volume density is
\(a(t)^m\) times round volume. The four pieces after the pulse form the
whole diagonal leg shown in Figure~\ref{fig:period}.

\subsection{The center-regular value equation}
Before analyzing the harmonic equation, we record one estimate
for matrix differences.  It does not require the matrices in a
product to commute.  All matrix norms in its applications are
Frobenius norms from the fixed round inner products; in fixed
finite dimensions they also control operator norms.

\begin{lemma}[A damped matrix difference]
\label{lem:damping}
Let \(E:[u,v]\to\mathbb R^{N\times N}\) be continuously
differentiable, and let \(F,\mathsf L,\mathsf R\) be continuous
matrix-valued functions on the same interval.  Suppose
\[
 E'=F-\mathsf L E-E\mathsf R
\]
there, where \(\mathsf L,\mathsf R\) are
symmetric and
\(\lambda_{\min}(\mathsf L)+\lambda_{\min}(\mathsf R)\geq\gamma>0\).
Then, for \(u\leq t\leq v\),
\begin{align}
 \|E(t)\|&\leq e^{-\gamma(t-u)}\|E(u)\|
       +\int_u^t e^{-\gamma(t-w)}\|F(w)\|\,dw,
       \label{eq:damping-pointwise}\\
 \int_u^v\|E(t)\|\,dt
 &\leq\gamma^{-1}
       \left(\|E(u)\|+\int_u^v\|F(t)\|\,dt\right).
       \label{eq:damping-integrated}
\end{align}
\end{lemma}

\begin{proof}
The Frobenius inner product gives
\[
 \frac12\frac{d}{dt}\|E\|^2
 =\langle E,F\rangle-\langle E,\mathsf L E\rangle
                         -\langle E,E\mathsf R\rangle
 \leq \|E\|\,\|F\|-\gamma\|E\|^2.
\]
The corresponding upper differential inequality for \(\|E\|\)
holds also at its zeros, by continuity (or by first replacing the
norm by \((\|E\|^2+\varepsilon^2)^{1/2}\) and letting
\(\varepsilon\downarrow0\)).  Multiplication by \(e^{\gamma t}\)
and integration gives \eqref{eq:damping-pointwise}.  Integrating
that estimate in \(t\) and interchanging the nonnegative integrals
gives \eqref{eq:damping-integrated}.
\end{proof}

For a harmonic function whose angular part belongs to \(V_l\), write
\[
 u(r,\omega)=\sum_{i=1}^{N_l}Y_i(\log r)e_{l,i}(\omega).
\]
We identify its coefficient vector \(Y(t)\) with
\(\sum_iY_i(t)e_{l,i}\in V_l\).
The function \(u\) is called \emph{center-regular} if it extends smoothly across
the origin.  The volume formula in Proposition~\ref{prop:metric} is the
reason that a closed equation holds on this fixed coefficient
space.  With \(t=\log r\), it is
\begin{equation}
 Y''+p(t)Y'-A_l(t)Y=0,\qquad
 p(t)=m-1+mb(t),\qquad
 A_l(t)=A_l(a(t),q(t),J(t)).
 \label{eq:harmonic-ode}
\end{equation}
Indeed, the radial density is \(r^m a(\log r)^m\).  Its logarithmic
derivative contributes \(m+mb(t)\) in the radial equation, and the
change \(t=\log r\) subtracts \(1\), giving \(p(t)\).
The density is constant in the angular variable, and
\eqref{eq:angular-operator} preserves \(V_l\), so no additional
angular terms or other harmonic degrees occur.

\begin{lemma}[The center-regular value equation]
\label{lem:regular}
There are a threshold \(j_{\rm reg}\geq j_{\rm geom}\) and constants
\(0<c_P<C_P\) such that the following holds for every
\(j_0\geq j_{\rm reg}\), every sequence of parameters in \(K\),
and every \(2\leq l\leq L\).  There is a unique smooth symmetric
matrix \(P_l(t)\) satisfying
\begin{equation}
 P_l'=A_l-pP_l-P_l^2,\qquad P_l=lI\quad(t\leq2).
 \label{eq:riccati}
\end{equation}
It satisfies
\begin{equation}
 c_PI\leq P_l(t)\leq C_PI
 \qquad(t\in\mathbb R,\ 2\leq l\leq L).
 \label{eq:riccati-barriers}
\end{equation}
The solutions of \(Y'=P_lY\) are exactly the center-regular
solutions of \eqref{eq:harmonic-ode}.  Their value vector at
any one time can be prescribed arbitrarily and determines the
solution uniquely.  At each pulse start one also has
\begin{equation}
 \|P_l(t_j)-\theta_l I\|\leq Cj^{-2}
 \qquad(j\geq j_0),
 \label{eq:round-reset}
\end{equation}
including the first pulse.  The constants and the threshold are
uniform over all preceding choices of the parameters in \(K\).
\end{lemma}

\begin{proof}
The matrices \(A_l(t)\) are symmetric.  For sufficiently large
allowed \(j_0\), the values of \(a\) and \(q\) lie in fixed compact
positive intervals, and only finitely many \(V_l\) are involved.
Positivity of \eqref{eq:angular-operator} therefore gives constants
\[
 0<\lambda_-I\leq A_l(t)\leq\lambda_+I.
\]
The radial condition \(b\geq-1/4\), together with \(b\leq0\), gives
\[
 0<p_-:=3m/4-1\leq p(t)\leq p_+:=m-1.
\]
Choose \(c_P>0\) below all the initial values \(l\) so that
\(\lambda_--p_+c_P-c_P^2>0\), and choose \(C_P\) above all of
them so that \(\lambda_+-p_-C_P-C_P^2<0\).  We may decrease
\(c_P\) and increase \(C_P\) so that every \(\theta_l\) also lies
strictly between them.

The Riccati equation preserves symmetry.  At a first contact with
the lower barrier, a unit vector \(v\) with \(P_lv=c_Pv\) satisfies
\[
 v^{\mathsf T}P_l'v
 \geq\lambda_--p_+c_P-c_P^2>0.
\]
At a contact with the upper barrier the corresponding derivative
is at most \(\lambda_+-p_-C_P-C_P^2<0\).  These strict inward
derivatives prevent an eigenvalue from crossing either barrier.
The bounds also keep \(P_l\) in a compact set of matrices on every
finite time interval, so local existence for the ordinary
differential equation continues for all later times.  Uniqueness follows
from the same local ordinary differential equation theorem.  In the
Euclidean region, \(A_l=B_lI\) and \(p=m-1\), so \(P_l=lI\)
indeed satisfies the equation there.

Let \(\mathcal F_l(t)\) be the fundamental matrix of \(Y'=P_lY\)
with \(\mathcal F_l(2)=I\).  Its determinant is
\[
 \det\mathcal F_l(t)
 =\exp\left(\int_2^t\operatorname{tr}P_l(w)\,dw\right)>0.
\]
Moreover,
\[
 Y''=(P_l'+P_l^2)Y=A_lY-pY',
\]
so its columns solve \eqref{eq:harmonic-ode}.  On \(t\leq2\)
they are \(e^{l(t-2)}\) times constant vectors.  Conversely, the
two scalar radial factors for \eqref{eq:harmonic-ode} in this
region are \(e^{lt}\) and \(e^{-(l+m-1)t}\).  Smoothness at the
origin eliminates the second factor.  Thus every center-regular
solution is one of the solutions \(Y'=P_lY\).  Invertibility of
\(\mathcal F_l(t)\) gives the assertion about arbitrary value data.

It remains to prove the reset estimate.  The schedule gives
\begin{equation}
 t_j=\frac{j^8}{8}+O(j^7),\qquad
 \sup_{t\in[t_j,t_{j+1}]}\left|\frac{t}{t_j}-1\right|\longrightarrow0.
 \label{eq:time-asymptotic}
\end{equation}
Times on period \(j\) and on the final rest immediately preceding
it are therefore comparable to \(j^8\).  From
\eqref{eq:scale-tail},
\begin{equation}
 a(t)-a_*=O(j^{-2}),\qquad p(t)-p_\infty=O(j^{-10})
 \label{eq:period-tail}
\end{equation}
uniformly on these intervals.  On a round interval set
\(E_0=P_l-\theta_l I\).  Subtracting \eqref{eq:round-root} from
\eqref{eq:riccati} gives the exact difference equation
\[
 E_0'=f_l^{\mathrm{rd}}(t)I-(pI+P_l)E_0-E_0\theta_l I,\qquad
 f_l^{\mathrm{rd}}(t)=a(t)^{-2}B_l-p(t)\theta_l-\theta_l^2.
\]
Here \(f_l^{\mathrm{rd}}=O(j^{-2})\) by \eqref{eq:period-tail}.  The symmetric
coefficients \(pI+P_l\) and \(\theta_l I\) have a fixed positive
sum of lower eigenvalue bounds.  Lemma~\ref{lem:damping}, on a final
rest of length \(T_{j-1}\), gives
\[
 \|P_l(t_j)-\theta_lI\|
 \leq C e^{-\gamma T_{j-1}}+Cj^{-2}=O(j^{-2}).
\]
The discrepancy at the start of the rest is bounded by
\eqref{eq:riccati-barriers}, independently of the history.  For the
first pulse, the interval up to \(t_{j_0}\) is entirely round.
Apply the same estimate on
\([t_{j_0}/2,t_{j_0}]\) once \(j_0\) is large.  Its length tends
to infinity and all its times are comparable to \(j_0^8\), so the
same bound follows.  This proves \eqref{eq:round-reset} with the
claimed uniformity.
\end{proof}

\subsection{The response of one pulse}
\label{subsec:pulse}

The reset brings the regular radial derivative close to the
round value \(\theta_l\), but its \(O(j^{-2})\) error cannot simply
be multiplied by the pulse duration \(T_j=j^6\).  The next proof
uses an exact scalar radial solution for the long common growth
and integrates the remaining, damped matrix error.  This is the
point at which the denominator \(p_\infty+2\theta_l\) in
\eqref{eq:round-root} enters.  Contraction properties of positive
matrix Riccati flows are classical; see Lawson--Lim~\cite[Theorem~8.5]{LawsonLim}.
We prove the forced estimates, including their parameter
derivatives, in the form used here.

Fix a period \(j\), hold all parameters in earlier periods fixed,
and let the current parameter \(h\) vary in \(K\).  Write
\[
 T=T_j,\qquad \tau=(t-t_j)/T,\qquad
 \eta_j=j^{-2}+T^{-1}.
\]
Let \(\mathcal B_{l,j}(h)\) be the value transmission for
\(Y'=P_lY\) from the start to the end of the pulse, and write
\(P_{l,e}(h)=P_l(t_j+T,h)\).  Define the scalar solution
\begin{equation}
 v_l'=a(t)^{-2}B_l-p(t)v_l-v_l^2,\qquad
 v_l(t_j)=\theta_l,\qquad
 \varphi_{l,j}=\int_{t_j}^{t_j+T}v_l(t)\,dt.
 \label{eq:scalar-baseline}
\end{equation}
This solution is independent of \(h\) and of the earlier
parameters.

\begin{lemma}[Pulse limit with one parameter derivative]
\label{lem:pulse}
For the fixed compact ball and pulse family,
\begin{equation}
 \begin{aligned}
 e^{-\varphi_{l,j}}\mathcal B_{l,j}(h)
     &=H_l(1,h)+O_{C^1}(\eta_j),\\
 P_{l,e}(h)&=\theta_lI+O_{C^1}(\eta_j).
 \end{aligned}
 \label{eq:pulse-limit}
\end{equation}
Here \(O_{C^1}(\eta_j)\) means that the norm and the norm of its
first derivative in the current parameter \(h\), uniformly on
\(K\), are at most \(C\eta_j\).  The constants and the large-\(j\)
threshold are independent of every admitted preceding history.
\end{lemma}

\begin{proof}
The scalar barriers used in Lemma~\ref{lem:regular} also keep \(v_l\)
between fixed positive constants.  Subtracting
\eqref{eq:round-root} from \eqref{eq:scalar-baseline} gives
\[
 (v_l-\theta_l)'=
 \bigl(a(t)^{-2}B_l-p(t)\theta_l-\theta_l^2\bigr)
 -(p+v_l+\theta_l)(v_l-\theta_l).
\]
Its initial value is zero.  The forcing is \(O(j^{-2})\) by
\eqref{eq:period-tail}, and the damping coefficient is uniformly
positive.  Thus
\begin{equation}
 v_l-\theta_l=O(j^{-2})
 \quad\text{throughout the pulse}.
 \label{eq:baseline-close}
\end{equation}
We retain \(v_l\) exactly in the value equation: even the small
pointwise difference in \eqref{eq:baseline-close} can have a large
integral over \(T_j\).

For the angular operator, expand \eqref{eq:angular-operator} at
\(q=1\).  Its first derivative in \(q\) is
\[
 a^{-2}\left(\frac{B_l}{m}I+D_J^2|_{V_l}\right).
\]
In particular the trace-free part has the positive coefficient
\(a^{-2}(D_J^2)_{\mathrm{tf}}\) that appears in
\eqref{eq:control-generators}.  Since \(q=1+z/T\) on the pulse and
\(a(t)-a_*=O(j^{-2})\), Taylor's formula gives
\begin{equation}
 \begin{aligned}
 A_l(t,h)
 &=a(t)^{-2}B_lI+\frac{p_\infty+2\theta_l}{T}R_l(\tau,h)
       +\mathcal E_{A,l}(t,h),\\
 \|\mathcal E_{A,l}\|_{C^1_h}
 &\leq C\left(\frac{j^{-2}}{T}+\frac1{T^2}\right).
 \end{aligned}
 \label{eq:angular-expansion}
\end{equation}
At most one bump is nonzero at any time.  In the gaps the angular
operator is exactly scalar.  Smooth boundedness of the profiles
and their first \(h\)-derivatives proves the uniform remainder
bound, also across those gaps.

Set
\[
 Q_l(t,h)=v_l(t)I+\frac1T R_l(\tau,h).
\]
It is symmetric, and it has a fixed positive lower eigenvalue
bound when \(j\) is large.  Its Riccati residual is
\[
 \rho_l=A_l-pQ_l-Q_l^2-Q_l'.
\]
The scalar terms in this expression cancel by
\eqref{eq:scalar-baseline}.  Using
\(\partial_t(R_l/T)=\partial_\tau R_l/T^2\), the remaining
identity is
\[
 \rho_l
 =-\frac{p-p_\infty+2(v_l-\theta_l)}{T}R_l
   -\frac{R_l^2+\partial_\tau R_l}{T^2}
   +\mathcal E_{A,l}.
\]
Equations \eqref{eq:baseline-close} and
\eqref{eq:angular-expansion} therefore imply
\begin{equation}
 \sup_{t,h}
 \bigl(\|\rho_l\|+\|D_h\rho_l\|\bigr)
 \leq\frac{C\eta_j}{T}.
 \label{eq:residual}
\end{equation}

Let \(E_l=P_l-Q_l\).  No matrices are commuted in the exact
difference equation
\begin{equation}
 E_l'=\rho_l-(pI+P_l)E_l-E_lQ_l.
 \label{eq:pulse-error-equation}
\end{equation}
The lower eigenvalue bounds for \(P_l,Q_l\), and \(p\) allow
Lemma~\ref{lem:damping} with a fixed damping rate.  The profiles
vanish near \(\tau=0\), so
\[
 E_l(t_j)=P_l(t_j)-\theta_lI=O(j^{-2})
\]
by \eqref{eq:round-reset}.  From
\eqref{eq:damping-pointwise}, \eqref{eq:damping-integrated}, and
\eqref{eq:residual}, we obtain
\begin{equation}
 \sup_{t,h}\|E_l(t,h)\|\leq C\eta_j,\qquad
 \sup_h\int_{t_j}^{t_j+T}\|E_l(t,h)\|\,dt\leq C\eta_j.
 \label{eq:error-integrated}
\end{equation}
The initial discrepancy contributes its size to this integral,
because it is damped; it does not contribute its size times \(T\).

We next control the derivative in a current parameter.  Derivatives
are taken with the earlier history fixed.  Write
\(E_{l,h}=D_hE_l\) and \(Q_{l,h}=D_hQ_l\); the calculation applies
to every unit direction in the finite parameter space.
The coefficients of the Riccati equation depend smoothly on \(h\)
on the pulse, so the ordinary differential equation has smooth
parameter dependence there; the uniform barriers keep its solutions
inside one fixed bounded set.
Differentiating \eqref{eq:pulse-error-equation} and using
\(D_hP_l=E_{l,h}+Q_{l,h}\) gives the exact regrouping
\begin{equation}
 E_{l,h}'=D_h\rho_l
 -(pI+P_l)E_{l,h}-E_{l,h}P_l
 -Q_{l,h}E_l-E_lQ_{l,h}.
 \label{eq:parameter-error-equation}
\end{equation}
Its initial value is zero: the past is fixed and
\(R_l(0,h)=D_hR_l(0,h)=0\).  Also
\(\|Q_{l,h}\|\leq C/T\).  Apply Lemma~\ref{lem:damping} to
\eqref{eq:parameter-error-equation}.  Its forcing has pointwise
norm at most \(C\eta_j/T\), by \eqref{eq:residual} and the first
bound in \eqref{eq:error-integrated}.  Its integrated norm is at
most
\[
 C\eta_j+\frac{C}{T}
   \int_{t_j}^{t_j+T}\|E_l\|\,dt
 \leq C\eta_j.
\]
Consequently
\begin{equation}
 \sup_{t,h}\|E_{l,h}\|\leq\frac{C\eta_j}{T},\qquad
 \sup_h\int_{t_j}^{t_j+T}\|E_{l,h}\|\,dt\leq C\eta_j.
 \label{eq:parameter-error-integrated}
\end{equation}
This argument supplies the derivative estimate directly, without
an a priori estimate for \(D_hP_l\).

To pass from coefficient errors to value transmissions, let
\(\mathcal B_l(t,h)\) be the pulse fundamental matrix up to time
\(t\), with value \(I\) at \(t_j\), and set
\[
 F_l(t,h)=
 e^{-\int_{t_j}^t v_l(w)\,dw}\mathcal B_l(t,h),\qquad
 G_l(t,h)=H_l((t-t_j)/T,h).
\]
Then
\[
 F_l'=(R_l/T+E_l)F_l,\qquad
 G_l'=(R_l/T)G_l,\qquad F_l(t_j)=G_l(t_j)=I.
\]
The integrated norms of \(R_l/T\) and \(D_hR_l/T\) are uniformly
bounded; those of \(E_l\) and \(E_{l,h}\) are \(O(\eta_j)\).
The elementary integral inequality
\(u(t)\leq a+\int_{t_j}^t b(w)u(w)\,dw\), \(b\geq0\), implies
\(u(t)\leq a\exp(\int_{t_j}^t b)\): the right-hand side of the
first inequality has derivative at most \(b\) times itself.
Applying this inequality to the fundamental-matrix equations and
their parameter derivatives bounds
\(F_l,G_l,D_hF_l,D_hG_l\) uniformly.

For completeness, \(Z_l=F_l-G_l\) satisfies
\[
 Z_l'=(R_l/T+E_l)Z_l+E_lG_l,\qquad Z_l(t_j)=0.
\]
The integral bound for \(E_l\) gives
\(\sup_t\|Z_l\|\leq C\eta_j\).  Differentiating this equation
introduces the additional terms
\[
 (D_hR_l/T+E_{l,h})Z_l+E_{l,h}G_l+E_lD_hG_l.
\]
Their integrated norms are \(O(\eta_j)\), using the preceding
uniform bounds and \eqref{eq:parameter-error-integrated}.
The same integral inequality gives
\(\sup_t\|D_hZ_l\|\leq C\eta_j\).  Evaluation at the pulse end
proves the first assertion of \eqref{eq:pulse-limit}.
At that end \(R_l(1,h)=D_hR_l(1,h)=0\).  Combining
\(P_l=Q_l+E_l\), \eqref{eq:baseline-close},
\eqref{eq:error-integrated}, and
\eqref{eq:parameter-error-integrated} proves the second assertion.
Every estimate used the earlier history only through the uniform
reset and barriers of Lemma~\ref{lem:regular}, proving the stated
uniformity.
\end{proof}

\subsection{Removing the full diagonal leg}
\label{subsec:leg}

The pulse estimate describes values and radial derivatives at
the pulse end.  The radial derivative need not yet equal
\(\theta_l\) times the value, so diagonal angular coefficients
on the next leg do not by themselves give a diagonal map on
these data.  We now factor the full leg exactly.  This separates
its possibly large unequal growth from the small error in the
entrance derivative.

Let \(\mathcal T_{l,j}(h)\) denote the value transmission of
center-regular solutions from \(t_j\) to \(t_{j+1}\).  On the
diagonal leg let \(U_{l,j}\) and \(V_{l,j}\) be the value
propagation matrices from its start to its end for initial
Cauchy data \((Y,Y')=(I,0)\) and \((0,I)\), respectively.  These
matrices are diagonal in the fixed dwell basis.  Define
\begin{equation}
 D_{l,j}=U_{l,j}+\theta_lV_{l,j},\qquad
 K_{l,j}=D_{l,j}^{-1}V_{l,j}.
 \label{eq:diagonal-reference}
\end{equation}
Thus \(D_{l,j}\) is the value propagation on the entire leg for
the reference data \((I,\theta_l I)\).

\begin{lemma}[Exact normalization of the full leg and dwell growth]
\label{lem:leg}
The matrices \(U_{l,j}\), \(V_{l,j}\), and \(D_{l,j}\) have
positive diagonal entries, and
\begin{equation}
 0\leq K_{l,j}\leq\theta_l^{-1}I.
 \label{eq:K-bound}
\end{equation}
They depend only on the prescribed diagonal leg, not on its
current or earlier pulse parameters.  The exact transmission
identity is
\begin{equation}
 D_{l,j}^{-1}\mathcal T_{l,j}(h)
 =\left[I+K_{l,j}\bigl(P_{l,e}(h)-\theta_lI\bigr)\right]
       \mathcal B_{l,j}(h).
 \label{eq:full-normalization}
\end{equation}
Moreover, uniformly for the finite set of \(l,i\),
\begin{equation}
 \log (D_{l,j})_{ii}
 =L_jd_{l,i}+O(1+L_jj^{-2}+T_j)
 =L_j\bigl(d_{l,i}+o(1)\bigr).
 \label{eq:diagonal-growth}
\end{equation}
\end{lemma}

\begin{proof}
Every operator on the leg is a function of \(I\) and
\(D_{J_0}^2\), and hence is diagonal in the fixed dwell
eigenbasis.  A scalar coordinate satisfies
\[
 y''+p(t)y'-A_i(t)y=0,\qquad A_i(t)>0.
\]
Multiplying the equation for \(y'\) by the positive integrating
factor \(e^{\int p}\) gives
\[
 \bigl(e^{\int p}y'\bigr)'=e^{\int p}A_i y.
\]
Starting with \((y,y')=(1,0)\), positivity of \(y\) makes \(y'\)
nonnegative until any supposed first zero of \(y\); this prevents
that zero.  Starting with \((0,1)\), the solution is positive
for small positive time and the same argument keeps it positive.
Thus the final entries of \(U_{l,j}\) and \(V_{l,j}\) are
positive.  Formula \eqref{eq:diagonal-reference} gives positivity
of \(D_{l,j}\), and entry by entry
\[
 0<\frac{(V_{l,j})_{ii}}
          {(U_{l,j})_{ii}+\theta_l(V_{l,j})_{ii}}
 \leq\theta_l^{-1}.
\]
This proves \eqref{eq:K-bound}.

For an incoming value vector \(y\) at the start of the period,
the data at the pulse end are
\[
 \bigl(\mathcal B_{l,j}y,\,
       P_{l,e}\mathcal B_{l,j}y\bigr).
\]
Their value at the end of the leg is
\[
 \mathcal T_{l,j}y
 =\bigl(U_{l,j}+V_{l,j}P_{l,e}\bigr)\mathcal B_{l,j}y.
\]
Since \(D_{l,j}=U_{l,j}+\theta_lV_{l,j}\), left multiplication
by \(D_{l,j}^{-1}\) gives \eqref{eq:full-normalization}.
The order \(K_{l,j}(P_{l,e}-\theta_lI)\) is part of this
identity; no commutation with \(P_{l,e}\) is used.  The large
matrix \(D_{l,j}\) has disappeared from the error multiplier,
leaving only the uniform bound \eqref{eq:K-bound}.

To estimate \(D_{l,j}\), take its \(i\)-th scalar solution with
initial data \((1,\theta_l)\).  The positivity just proved
allows its logarithmic derivative \(w=y'/y\) throughout the
leg.  It satisfies
\[
 w'=A_{l,i}(t)-p(t)w-w^2.
\]
The scalar barrier argument in Lemma~\ref{lem:regular}, with
initial value \(\theta_l\), keeps \(w\) between fixed positive
constants.  During the dwell,
\[
 A_{l,i}(t)=\lambda_{l,i}+O(j^{-2}),\qquad
 \lambda_{l,i}=d_{l,i}(d_{l,i}+p_\infty),
\]
by \eqref{eq:period-tail} and the fixed value \(q=q_*\).
For \(e=w-d_{l,i}\), subtraction gives
\[
 e'=F_{l,i}(t)-(p+w+d_{l,i})e,\qquad
 F_{l,i}=A_{l,i}-\lambda_{l,i}
                 -(p-p_\infty)d_{l,i}=O(j^{-2}).
\]
The value of \(e\) at the dwell entrance is bounded, and its
damping coefficient has a fixed positive lower bound.  The
scalar version of \eqref{eq:damping-integrated} yields
\[
 \int_{\text{dwell }j}|w-d_{l,i}|\,dt
 \leq C+C L_jj^{-2}.
\]
The two ramps and the final rest have total duration \(3T_j\),
and their contribution to \(\int w\) is \(O(T_j)\).  Since the
solution starts at \(1\) and ends at \((D_{l,j})_{ii}\),
integration of \(w=(\log y)'\) gives the first equality in
\eqref{eq:diagonal-growth}.  The second follows from
\(T_j/L_j=j^{-1}\to0\) and \(L_j\to\infty\).  All constants
are uniform over the fixed finite list of coordinates.
\end{proof}

\subsection{Tuning every complete period exactly}
\label{subsec:tuning}

The preceding identity puts the actual full-period map in the
same local coordinate problem as the ideal word.  Its remaining
error tends to zero with one parameter derivative, uniformly
over the past.  This allows the target cycle to be attained
successively in every period.

\begin{proposition}[Exact full-period transmission]
\label{prop:transmission}
For finite spectral data as specified at the end of Section~\ref{sec:counting},
choose the word and compact pulse family above. There are a starting
index \(j_0\) and parameters \(h_j\in K\), for every \(j\geq j_0\),
such that the center-regular value transmissions of the resulting metric
satisfy
\begin{equation}
 \mathcal T_{l,j}=s_{l,j}D_{l,j}\Pi_l,\qquad
 s_{l,j}>0,\qquad |\log s_{l,j}|\leq C T_j
 \quad(2\leq l\leq L,\ j\geq j_0).
 \label{eq:exact-transmission}
\end{equation}
The positive diagonal matrices \(D_{l,j}\) propagate the reference
Cauchy data \((I,\theta_l I)\) over the entire leg following the pulse;
they satisfy \eqref{eq:diagonal-growth}. The regular value equation has the
bounds \eqref{eq:riccati-barriers}. All estimates used to choose \(j_0\)
are uniform over preceding admitted controls for this fixed finite family.
\end{proposition}

\begin{proof}
First fix an admitted past and vary the current parameter \(h\).
By \eqref{eq:full-normalization},
\[
 e^{-\varphi_{l,j}}D_{l,j}^{-1}\mathcal T_{l,j}(h)
 =\left[I+K_{l,j}\bigl(P_{l,e}(h)-\theta_lI\bigr)\right]
       e^{-\varphi_{l,j}}\mathcal B_{l,j}(h).
\]
The bound \eqref{eq:K-bound} and Lemma~\ref{lem:pulse} imply
\begin{equation}
 e^{-\varphi_{l,j}}D_{l,j}^{-1}\mathcal T_{l,j}(h)
 =H_l(1,h)+O_{C^1}(\eta_j).
 \label{eq:full-period-limit}
\end{equation}
The estimate is uniform in the preceding history.  Each
\(\mathcal T_{l,j}\) is a value fundamental matrix of
\(Y'=P_lY\), so it has positive determinant; so does \(D_{l,j}\).
Determinant normalization is therefore defined for their
quotient.  The matrices \(H_l(1,h)\), \(h\in K\), form a compact
subset of the positive-determinant matrices.  On a neighborhood
of this compact set, \(\mathcal N_l\) is smooth with bounded
derivatives.  Normalization is unchanged by the positive scalar
\(e^{-\varphi_{l,j}}\), so \eqref{eq:full-period-limit} gives
\begin{equation}
 \Psi_j(h):=
 \bigl(\mathcal N_l(D_{l,j}^{-1}\mathcal T_{l,j}(h))\bigr)_{l=2}^{L}
 =\Phi(h)+O_{C^1}(\eta_j)
 \quad\text{in }\mathcal G.
 \label{eq:group-limit}
\end{equation}
The assertion is first an estimate for the ambient matrix
entries and their derivatives, and hence also in any fixed
smooth chart containing the compact image under consideration.

We spell out the uniform inversion.  Use Euclidean norms on the
parameter and coordinate spaces, and the associated operator norms
for their derivatives.  Choose a coordinate chart
\(\zeta\) near \(\Pi\) in \(\mathcal G\) with \(\zeta(\Pi)=0\),
and let \(f=\zeta\circ\Phi\) and \(A_0=Df(0)\).  By
\eqref{eq:ideal-word}, \(A_0\) is invertible.  Choose a closed
ball \(\overline B_r\) centered at \(0\), contained in the
interior of \(K\), whose \(\Phi\)-image lies compactly inside
the chart and for which
\[
 \sup_{h\in\overline B_r}
 \|I-A_0^{-1}Df(h)\|\leq\frac14.
\]
For every sufficiently large \(j\), uniformly in the admitted
past, \(f_j=\zeta\circ\Psi_j\) is defined on this ball and
\[
 \sup_{h\in\overline B_r}
 \|A_0^{-1}(Df_j(h)-Df(h))\|\leq\frac14,\qquad
 \|A_0^{-1}f_j(0)\|\leq\frac r2.
\]
This follows from \eqref{eq:group-limit} and \(\eta_j\to0\).
The map
\[
 \mathcal C_j(h)=h-A_0^{-1}f_j(h)
\]
has derivative of norm at most \(1/2\) on the ball.  Also
\[
 \|\mathcal C_j(h)\|
 \leq\|\mathcal C_j(0)\|+\tfrac12\|h\|
 \leq r
 \qquad(h\in\overline B_r).
\]
It is therefore a contraction of the closed ball into itself.
Its iterates converge to a fixed point \(h_j\), for which
\(f_j(h_j)=0\), or \(\Psi_j(h_j)=\Pi\).

All the preceding estimates were uniform over the past.
Choose \(j_0\) once, large enough for them and for the geometric
input, after the fixed word and parameter ball have been chosen.
The history before \(j_0\) is round.  Choose \(h_{j_0}\) by the
contraction above, and repeat after each finite chosen history.
This inductively defines an infinite admitted sequence and
hence a single smooth metric of Proposition~\ref{prop:metric}.
Undoing determinant normalization at each fixed point gives
\[
 D_{l,j}^{-1}\mathcal T_{l,j}=s_{l,j}\Pi_l,\qquad
 s_{l,j}=\det(D_{l,j}^{-1}\mathcal T_{l,j})^{1/N_l}>0.
\]
By \eqref{eq:full-period-limit}, the determinant of
\(e^{-\varphi_{l,j}}D_{l,j}^{-1}\mathcal T_{l,j}\) stays
bounded above and away from zero on the fixed compact family.
Thus
\[
 \log s_{l,j}=\varphi_{l,j}+O(1)=O(T_j),
\]
because \(v_l\) is uniformly bounded.  This proves
\eqref{eq:exact-transmission}.  The equality is for the complete
period; no residual change of direction is passed to a later
period.

\end{proof}
